% !TEX root = main.tex
\subsection{Syntax}\label{subsection:sms-syntax}

We adopt the partial-dynamic model~\cite{KiayiasY06},
where users may join at any time but revocation is not supported;
a fully dynamic model would additionally allow membership expiry.

\medskip
\noindent\textbf{Parties.}
The \emph{function-issuing authority} $\FA$ holds master function key $\mfk$
and enrolls users by certifying each user's attribute and hidden policy.
The \emph{tracing authority} $\TA$ holds tracing key $\tsk$
and recovers signer identity or role from an opening tag when authorized.
Each \emph{user} holds a certified attribute--policy pair $(\attr, F)$
and chooses a public ad hoc function $G$ at each signing.
\emph{Verifiers} check signatures publicly using $\pp$ and the declared $G$.

\medskip
\noindent\textbf{Signing process.}
To sign a message $M \in \MessageSpace$ with auxiliary witness
$w \in \WitnessSpace$ supporting the computations of $F$ and $G$,
the signer computes
\[
(\tM, \op) = P\bigl(F(M, \attr, w),\, G(M, \attr, w)\bigr),
\]
where $P$ is the public combiner fixed in $\pp$
and $\tM \in \TMessageSpace$ is the derived message
(equal to $M$ when $P$ returns the input message unchanged).
If $\op \in \OpeningSetSignable$ the signer issues a signature $\Sigma$ on $(\tM, G)$,
which also serves as a zero-knowledge proof of correct evaluation;
if $\op \notin \OpeningSetSignable$ signing is blocked.
The attribute $\attr$ and policy $F$ remain hidden in $\Sigma$;
only $\TA$ learns $\op$ via $\SMSOpen$.

\begin{definition}[\SMS]\label{definition:SMS-syntax}
An \SMS scheme is defined by the algorithms below.
%a tuple $\SMSScheme = (\SMSSetup, \SMSJoin, \SMSSign, \SMSVerify, \SMSOpen)$ with the following algorithms.

\begin{description}

\item[$\SMSSetup(1^\lambda, N)$:]
On input security parameter $1^\lambda$ and maximum user count $N$,
run $(\vk, \mfk) \gets \SKeyGen(1^\lambda)$ for $\FA$
and $(\ek, \tsk) \gets \EKeyGen(1^\lambda)$ for $\TA$,
initialize $\State \coloneqq \emptyset$, and publish the public parameter
\[
\pp = \bigl(\gpk,\; \MessageSpace,\; \TMessageSpace,\; \AttrSpace,\;
\WitnessSpace,\; \FFamily,\; \GFamily,\; P,\;
\OpeningSet,\; \OpeningSetSignable\bigr),
\]
where $\gpk$ incorporates $\vk$ and $\ek$,
$\MessageSpace$ and $\TMessageSpace$ are the input and derived message spaces,
$\AttrSpace$ and $\WitnessSpace$ are the attribute and witness spaces,
$\FFamily$ and $\GFamily$ are the hidden and ad hoc function families,
$P$ is the public combiner,
$\OpeningSet$ is the set of all opening values,
and $\OpeningSetSignable \subsetneq \OpeningSet$ is the signable subset.
Output $(\pp, \State, \mfk, \tsk)$.

The strict inclusion ensures that $\OpeningSet \setminus \OpeningSetSignable$ is non-empty,
so the combiner $P$ can block signing by outputting a value outside $\OpeningSetSignable$,
encoding a policy rejection without producing a signature.

The unified $\SMSSetup$ is retained here for notational convenience.
In deployments where $\FA$ and $\TA$ are distinct,
$\SMSSetup$ can be split into two independent key generation algorithms:
$\SMSSetup_\FA(1^\lambda)$ outputting $(\vk, \mfk)$
and $\SMSSetup_\TA(1^\lambda)$ outputting $(\ek, \tsk)$,
with $\pp$ assembled from their public components,
such that the two authority keys $\mfk$ and $\tsk$ are held by separate parties $\FA$ and~$\TA$.

\item[$\SMSJoin\interaction{\User(1^\lambda)}{\FA(\mfk)}(\pp, \State, \attr, F)$:]
$\FA$ enrolls user $\User$ with attribute $\attr \in \AttrSpace$
and assigns function $F \in \FFamily$.
Upon successful completion:
\begin{enumerate}
\item $\User$ obtains $(\secattr, \certattr)$,
where $\secattr$ is the user secret key and
$\certattr$ is a certificate binding $(\attr, F)$;
\item $\FA$ updates $\State' \coloneqq (\State \| (\attr, F, \transcript))$, 
where $\transcript$ is the transcript.
The updated state $\State'$ replaces $\State$ in subsequent executions.
\end{enumerate}
If $\SMSJoin$ fails, $\User$ receives $\perp$, and $\State$ is unchanged.
By convention, $\certattr$ encodes $(\attr, F)$ in the clear so that
any party may read $(\attr, F)$ from $\certattr$;
the pair $(\attr, F)$ need only be concealed from signatures, not from certificates.
We write $\secattr \leftrightharpoons_\pp \certattr$ to indicate that
$(\secattr, \certattr)$ was produced by an honest execution of $\SMSJoin$
under $\pp$, 
$\secattr \not\leftrightharpoons_\pp \certattr$ otherwise.

%Every signature produced by $\User$ thereafter must prove correct evaluation under $\pp$ relative to a valid $\certattr$.

\item[$\SMSSign(\pp, \certattr, \secattr, G, M, w)$:]
Given $\pp$, a certificate $\certattr$ along with a secret $\secattr$,
a function $G \in \GFamily$, a message $M \in \MessageSpace$,
and a witness $w \in \WitnessSpace$,
output a derived message $\tM \in \TMessageSpace$ and a signature $\Sigma$.
Return $\perp$ if it fails.

\item[$\SMSVerify(\pp, G, \tM, \Sigma)$:]
Output $1$ if $\Sigma$ is a valid \SMS signature on $\tM \in \TMessageSpace$
under function $G \in \GFamily$ and parameters $\pp$;
output $0$ otherwise.

\item[$\SMSOpen(\pp, \tsk, G, \tM, \Sigma)$:]
Using tracing key $\tsk$, extract and return the opening value $\op \in \OpeningSet$
from a valid signature $\Sigma$ on derived message $\tM$ under function $G$.
The information recoverable from $\op$ ranges from a role or partial identity
to the full identity, depending on the configuration encoded in~$P$.

\end{description}
\end{definition}

\medskip
\noindent\textbf{Illustration.}
A central bank issues digital currency (CBDC) to verified citizens.
Each citizen's signing key encodes a certified hidden policy
$F(M, \attr, w)$ that checks transaction validity against her know-your-customer compliance tier $\attr$:
whether the transaction amount is within her spending limit,
the merchant category is permitted,
and the note has not been spent before (via witness $w$).
At each payment, the citizen additionally chooses a public function $G$
encoding the transaction's public validity conditions,
\eg, that the amount is positive and the merchant is registered.
The combiner $P$ outputs the transaction record $\tM$
and an opening tag $\op$ reflecting the transaction's risk level:
routine payments yield $\op = \mathsf{pseudonym}$,
exposing only a per-citizen unlinkable tag to the central bank;
transactions flagged by $F$ as anomalous yield $\op = \mathsf{id}$,
exposing the citizen's full identity to the regulator.
If $F$ rejects the transaction outright, $P$ outputs $\op \notin \OpeningSetSignable$,
blocking the payment without producing a signature.
The spending policy $F$ remains hidden from merchants and the public;
verifiers confirm only that $G$ approved the transaction
under some certified hidden policy.
%issued by the central bank.
